\section*{الفصل \ref{ch:g12:respiration-fermentation} --- التنفس الخلوي والتخمر}

\begin{solution}{exo:g12:respiration-fermentation:1}
\omterm{def:g10:universal-dna:nucleotide}{نوكليوتيد} له ثلاث مجموعات فوسفاتية على التوالي، وهو عملة \omterm{def:g10:cells-common-unit:cell}{الخلية}
للطاقة. واستعماله ينزع الفوسفات الأخير، فيحرر نحو \qty{30}{kJ} في
المول ليقود مسارًا، ويترك ADP ليعاد شحنه.
\end{solution}

\begin{solution}{exo:g12:respiration-fermentation:2}
\omterm{prop:g12:respiration-fermentation:glycolysis}{تحلل السكر} في \omterm{def:g10:cells-common-unit:cell}{السيتوبلازم}: 2 \omterm{def:g12:respiration-fermentation:atp}{ATP} و2 NAD مرجَع، وبيروفاتان. \omterm{prop:g12:respiration-fermentation:mitochondrion}{ودورة
كريبس} في المادة الأساسية \omterm{def:g10:cells-common-unit:organelle}{للمتقدرة}: ثنائي أكسيد كربون، وحاملان
محمَّلان، و2 \omterm{def:g12:respiration-fermentation:atp}{ATP}. والسلسلة التنفسية في الغشاء الداخلي: الحاملان
يفرغان على الأكسجين، وماء، ونحو 26 \omterm{def:g12:respiration-fermentation:atp}{ATP}.
\end{solution}

\begin{solution}{exo:g12:respiration-fermentation:3}
غشاءان، الداخلي منهما مطوي إلى أعراف، يحيطان بالمادة الأساسية. المادة
الأساسية: \omterm{prop:g12:respiration-fermentation:mitochondrion}{دورة كريبس}. والغشاء الداخلي: السلسلة التنفسية وتركيب \omterm{def:g12:respiration-fermentation:atp}{ATP}.
\end{solution}

\begin{solution}{exo:g12:respiration-fermentation:4}
هو يفرغ NAD المرجَع الآتي من \omterm{prop:g12:respiration-fermentation:glycolysis}{تحلل السكر} على البيروفات، ليستطيع \omterm{prop:g12:respiration-fermentation:glycolysis}{تحلل
السكر} أن يستمر بلا أكسجين. والمردود: جزيئا \omterm{def:g12:respiration-fermentation:atp}{ATP} الآتيان من \omterm{prop:g12:respiration-fermentation:glycolysis}{تحلل السكر}
لا غير، مع لاكتات أو إيثانول و$\mathrm{CO_2}$.
\end{solution}

\begin{solution}{exo:g12:respiration-fermentation:5}
نحو 30 \omterm{def:g12:respiration-fermentation:atp}{ATP} في مقابل 2. والتنفس يترك ثنائي أكسيد الكربون والماء؛
\omterm{def:g10:cell-metabolism:fermentation}{والتخمر} يترك لاكتات، أو إيثانولًا وثنائي أكسيد كربون، ما زالت غنية
بالطاقة.
\end{solution}

\begin{solution}{exo:g12:respiration-fermentation:6}
الغلوكوز: لا أثر --- \omterm{def:g10:cells-common-unit:organelle}{فالمتقدرات} لا تستطيع استعماله. والبيروفات: يبدأ
استهلاك الأكسجين --- فهو \omterm{def:g11:enzymes-and-phenotype:enzyme}{ركيزة} \omterm{prop:g12:respiration-fermentation:mitochondrion}{المتقدرة}. وADP: يتسارع الاستهلاك ---
فاستعمال الأكسجين مقترن بتركيب \omterm{def:g12:respiration-fermentation:atp}{ATP}. والسيانيد: يتوقف الاستهلاك ---
فالسلسلة التنفسية هي موضع استعمال الأكسجين.
\end{solution}

\begin{solution}{exo:g12:respiration-fermentation:7}
\omterm{prop:g12:respiration-fermentation:glycolysis}{تحلل السكر}، وهو يحوّل الغلوكوز إلى بيروفات، يجري في \omterm{def:g10:cells-common-unit:cell}{السيتوبلازم}؛
\omterm{prop:g12:respiration-fermentation:mitochondrion}{والمتقدرة} تفتقر إلى إنزيماته وتأخذ البيروفات لا الغلوكوز.
\end{solution}

\begin{solution}{exo:g12:respiration-fermentation:8}
كل غلوكوز يحمّل جزيئي NAD؛ وللخلية مخزون صغير من NAD لا غير، فما إن
يحمَّل كله حتى يفقد \omterm{prop:g12:respiration-fermentation:glycolysis}{تحلل السكر} حاملًا يسلّمه هيدروجينه فيتوقف. \omterm{def:g10:cell-metabolism:fermentation}{والتخمر}
يفرغ NAD على البيروفات ويبقي السبيل يعمل.
\end{solution}

\begin{solution}{exo:g12:respiration-fermentation:9}
التنفس: $60/30 = 2$ غلوكوز في الثانية. \omterm{def:g10:cell-metabolism:fermentation}{والتخمر}: $60/2 = 30$ غلوكوزًا
في الثانية، أي أكثر بخمسة عشر ضعفًا.
\end{solution}

\begin{solution}{exo:g12:respiration-fermentation:10}
السيانيد يسد السلسلة التنفسية، فلا تبلغ إلكترونات الأكسجين ويضيع نحو
26 من الثلاثين \omterm{def:g12:respiration-fermentation:atp}{ATP} لكل غلوكوز في الحال؛ \omterm{def:g10:cell-metabolism:fermentation}{والتخمر} لا يستطيع سد النقص.
والدماغ والقلب يستعملان \omterm{def:g12:respiration-fermentation:atp}{ATP} أسرع من غيرهما ولهما أصغر احتياط، فيخفقان
في غضون دقائق.
\end{solution}

\begin{solution}{exo:g12:respiration-fermentation:11}
ثنائي أكسيد الكربون يتحرر في المادة الأساسية \omterm{def:g10:cells-common-unit:organelle}{للمتقدرة}، \omterm{prop:g12:respiration-fermentation:mitochondrion}{بدورة كريبس}
(وبالخطوة التي تدخل البيروفات إليها): ولا شيء منه في \omterm{prop:g12:respiration-fermentation:glycolysis}{تحلل السكر}.
والماء يصنع في نهاية السلسلة التنفسية، حين تنضم الإلكترونات
والبروتونات إلى الأكسجين.
\end{solution}

\begin{solution}{exo:g12:respiration-fermentation:12}
نالت العضلة 2 \omterm{def:g12:respiration-fermentation:atp}{ATP} لكل غلوكوز خمّرته؛ وإعادة بناء ذلك الغلوكوز من
اللاكتات تكلف الكبد نحو 6 \omterm{def:g12:respiration-fermentation:atp}{ATP}، تدفع بتنفس وقود آخر. فالجسم أنفق 6
لينال 2: أي إن \omterm{def:g10:cell-metabolism:fermentation}{التخمر} يستقرض طاقة لا بد أن ترد بفائدة.
\end{solution}

\begin{solution}{exo:g12:respiration-fermentation:13}
في الهواء يعطي التنفس 30 \omterm{def:g12:respiration-fermentation:atp}{ATP} لكل غلوكوز، فيلزم غلوكوز قليل ولا يحوَّل
شيء منه إلى إيثانول؛ وإذا حبست الخميرة أعطى \omterm{def:g10:cell-metabolism:fermentation}{التخمر} 2، فيستهلك غلوكوز
أكثر بخمسة عشر ضعفًا ويصنع الإيثانول. والعجينة تنتفخ بثنائي أكسيد
الكربون، وكلا الطريقين ينتجه (التنفس ستة لكل غلوكوز، \omterm{def:g10:cell-metabolism:fermentation}{والتخمر} اثنين)؛
وفي العجينة الكثيفة ينفد الأكسجين سريعًا على كل حال.
\end{solution}

\begin{solution}{exo:g12:respiration-fermentation:14}
حرارةً: فطاقة السلسلة، وقد كفّت عن أن تلتقط \omterm{def:g12:respiration-fermentation:atp}{ATP}، تتحرر مباشرة. والمولود
الجديد يفقد الحرارة سريعًا ولا يحسن الارتعاد؛ فالدهن البني مدفأة
مركَّبة فيه.
\end{solution}

\begin{solution}{exo:g12:respiration-fermentation:15}
الإلكترونات: \omterm{def:g10:cell-metabolism:autotrophy}{التركيب الضوئي} يأخذها من الماء ويضعها في السكر؛ والتنفس
يأخذها من السكر ويعطيها الأكسجين. والغازات: \omterm{def:g10:cell-metabolism:autotrophy}{التركيب الضوئي} يستهلك
$\mathrm{CO_2}$ ويحرر $\mathrm{O_2}$؛ والتنفس بالعكس. والطاقة:
\omterm{def:g10:cell-metabolism:autotrophy}{التركيب الضوئي} يدخل الضوء ويخزنه في السكر؛ والتنفس يخرجه من السكر
على صورة \omterm{def:g12:respiration-fermentation:atp}{ATP} وحرارة. والعضيتان: \omterm{def:g12:photosynthesis:chloroplast}{البلاستيدة الخضراء} \omterm{prop:g12:respiration-fermentation:mitochondrion}{والمتقدرة}. وكل
منهما يستهلك ما ينتجه الآخر: \omterm{def:g10:cell-metabolism:autotrophy}{فالتركيب الضوئي} كان سيستنفد
$\mathrm{CO_2}$ والتنفس كان سيستنفد $\mathrm{O_2}$ والغذاء.
\end{solution}

\begin{solution}{pb:g12:respiration-fermentation:1}
\textbf{1.} ثلث \qty{1000}{W} هو \qty{300}{W}: أي \qty{18}{kJ} في
الدقيقة، أي \qty{0.6}{mol} من \omterm{def:g12:respiration-fermentation:atp}{ATP} في الدقيقة.

\textbf{2.} نحو \qty{300}{g} في الدقيقة --- في مقابل غرامات قليلة
محمولة: فالمخزون يعاد تدويره تدويرًا متصلًا.

\textbf{3.} \qty{50}{g} هي \qty{0.1}{mol}: فيعاد شحن كل جزيء نحو 6
مرات في الدقيقة.

\textbf{4.} $0.6/30 = \qty{0.02}{mol}$ من الغلوكوز في الدقيقة (أي نحو
\qty{3.6}{g}).

\textbf{5.} $6 \times 0.02 = \qty{0.12}{mol}$ من $\mathrm{O_2}$، أي
نحو \qty{2.9}{L} في الدقيقة --- دون أقصاها البالغ \qty{4}{L/min}،
وموافق لقيمة \qty{20}{kJ} لكل لتر أكسجين في
\cref{ch:g10:exercise-and-energy}: فاستطاعة \qty{1000}{W} يمكن أن
يديمها التنفس.

\textbf{6.} 2 في \omterm{def:g10:cells-common-unit:cell}{السيتوبلازم}، و28 في \omterm{prop:g12:respiration-fermentation:mitochondrion}{المتقدرة}: أي نحو 93\% من \omterm{def:g12:respiration-fermentation:atp}{ATP}.

\textbf{7.} $6 \times 0.02 = \qty{0.12}{mol}$ من $\mathrm{CO_2}$ في
الدقيقة، كلها من \omterm{prop:g12:respiration-fermentation:mitochondrion}{دورة كريبس} ومن دخول البيروفات في المادة الأساسية.

\textbf{8.} $30 \times 30 = \qty{900}{kJ}$ من \qty{2870}{kJ}، أي نحو
الثلث؛ والباقي حرارة تطرحها بالتعرق وبجريان الدم إلى جلدها.

\textbf{9.} بروتينات السلسلة \omterm{def:g11:enzymes-and-phenotype:enzyme}{والإنزيمات} التي تصنع \omterm{def:g12:respiration-fermentation:atp}{ATP} تجلس في الغشاء؛
ومعدل تركيب \omterm{def:g12:respiration-fermentation:atp}{ATP} يتناسب مع المساحة التي تحملها.

\textbf{10.} يرفع المعدل الذي يستطيع به التنفس أن يصنع \omterm{def:g12:respiration-fermentation:atp}{ATP} (وحصة الدسم
التي تستطيع حرقها)، ومن ثم الاستطاعة التي تستطيع إدامتها بلا \omterm{def:g10:cell-metabolism:fermentation}{تخمر}؛
ولا يغير المردود لكل غلوكوز ولا الأكسجين لكل غلوكوز.

\textbf{11.} $800 \times 30 = \qty{24}{kJ}$: أي \qty{0.8}{mol} من \omterm{def:g12:respiration-fermentation:atp}{ATP}.

\textbf{12.} التنفس: $400 \times 30 = \qty{12}{kJ}$، أي
\qty{0.4}{mol}؛ ولا بد أن يمد \omterm{def:g10:cell-metabolism:fermentation}{التخمر} بالكمية \qty{0.4}{mol} الأخرى.

\textbf{13.} $0.4/2 = \qty{0.2}{mol}$ من الغلوكوز، ويترك \qty{0.4}{mol}
من اللاكتات.

\textbf{14.} التنفس: $1/30$ غلوكوز لكل \omterm{def:g12:respiration-fermentation:atp}{ATP}؛ \omterm{def:g10:cell-metabolism:fermentation}{والتخمر}: $1/2$ --- أي
غلوكوز أكثر بخمسة عشر ضعفًا. وتقبله العضلة لأن \omterm{def:g12:respiration-fermentation:atp}{ATP} مطلوب الآن
\omterm{def:g10:exercise-and-energy:glycogen}{والغليكوجين} حاضر؛ والدين يرد لاحقًا.

\textbf{15.} لا بد أن تؤكسد اللاكتات أو يعاد بناؤها غلوكوزًا، وذلك
يكلف \omterm{def:g12:respiration-fermentation:atp}{ATP} يصنعه التنفس: فيبقى استهلاك الأكسجين فوق مستوى الراحة حتى
يسدد الدين.

\textbf{16.} \qty{500}{W} مدة ساعة هي \qty{1800}{kJ}: أي $1800/38
\approx \qty{47}{g}$ من الدسم.

\textbf{17.} الدسم أكثر إرجاعًا --- أي أغنى بالهيدروجين --- فيحتاج كل
غرام منه أكسجينًا أكثر ليؤكسد أكسدة تامة؛ وهو يعطي، لكل لتر أكسجين،
طاقة أقل قليلًا مما يعطي الغلوكوز.

\textbf{18.} \omterm{prop:g12:respiration-fermentation:glycolysis}{تحلل السكر} وتخمره: فالدسم لا يخمَّر، ومن ثم لا تستطيع
سرعة نهائية، وهي تعمل \omterm{def:g10:cell-metabolism:fermentation}{بالتخمر}، أن تستعمل الدسم البتة --- بل الغلوكوز
\omterm{def:g10:exercise-and-energy:glycogen}{والغليكوجين} لا غير.

\textbf{19.} يحرر الكبد الغلوكوز من غليكوجينه ويعيد بناء بعضه من
اللاكتات، فيبقي غلوكوز الدم قرب \qty{1}{g/L}
(\cref{ch:g12:glucose-and-diabetes})؛ فيحمى إمداد الدماغ.

\textbf{20.} نحو \qty{0.6}{mol} من \omterm{def:g12:respiration-fermentation:atp}{ATP} في الدقيقة؛ ونحو \qty{3}{L} من
الأكسجين في الدقيقة تدفع ثمنها؛ والسلسلة التنفسية في الغشاء الداخلي
\omterm{def:g10:cells-common-unit:organelle}{للمتقدرة} تصنع نحو تسعة أعشارها.
\end{solution}
